% RESULTADO_RECUPERADO desde II REV09. Véase MAPA_PROCEDENCIA.json.
% Selección de líneas y cambios mecánicos explícitos; ninguna prueba nueva.
\section{Operador de área y reconstrucción de la información sectorial}
\label{v:ant:sec:area}
El reagrupamiento de seis trits construye una biyección entre
$(\Z/9\Z)^3$ y $(\Z/27\Z)^2$, ambos de cardinal $729$.
La biyección no se presenta como isomorfismo aditivo. Sobre un bloque
de borde $9\times9$, la fuente descompone 36 enlaces en dos
conjuntos de 18, etiquetados por los canales $90$ y $120$.

\subsection{Reagrupamiento trítico e incidencia del borde}
\label{v:ant:sec:ii-area-incidencia-explicita}

La correspondencia conserva una palabra ordenada, además de su cardinal.
Escribamos cada coordenada \(x_j\in\mathbb Z/9\mathbb Z\), para
\(j=1,2,3\), mediante sus representantes posicionales únicos
\(x_j=r_{2j-1}+3r_{2j}\), con \(r_k\in\{0,1,2\}\). El mapa es
\begin{equation}
 \beta_{729}(x_1,x_2,x_3)
 =\bigl(r_1+3r_2+9r_3,\ r_4+3r_5+9r_6\bigr)
 \in(\mathbb Z/27\mathbb Z)^2.
 \label{v:ant:eq:ii-area-beta729}
\end{equation}

\begin{proposition}[Correspondencia posicional entre interior y borde]
\label{v:ant:prop:ii-area-beta729}
La aplicación \(\beta_{729}\) es biyectiva y conserva el orden de los
seis trits. Su inversa extrae las dos ternas de dígitos y las reagrupa
en tres parejas.
\end{proposition}
\begin{proof}
La expansión posicional de los representantes \(0,\ldots,8\) y
\(0,\ldots,26\) es única. Extraer los seis dígitos, reagruparlos y
extraerlos de nuevo devuelve la misma palabra. La operación inversa
restituye así cada una de las tres coordenadas iniciales.
\end{proof}

El acarreo conserva su posición en la palabra enriquecida. Las operaciones
aditivas de los dos empaquetados tienen fronteras diferentes: por ejemplo,
\(\beta_{729}(0,2,0)=(18,0)\) y
\(\beta_{729}(0,1,0)=(9,0)\), cuya suma en el toro de orden
veintisiete es \((0,0)\), mientras
\(\beta_{729}(0,3,0)=(0,1)\). Esta distinción permite reutilizar
la correspondencia posicional manteniendo tipadas las operaciones.

En el toro \((\mathbb Z/27\mathbb Z)^2\), el bloque
\(B_\partial=\{0,\ldots,8\}^2\) tiene los conjuntos de enlaces
orientados
\begin{align}
 \partial B_{90}
 &=\{((-1,j),(0,j)),\ ((8,j),(9,j)):0\le j\le8\},\\
 \partial B_{120}
 &=\{((i,-1),(i,0)),\ ((i,8),(i,9)):0\le i\le8\}.
 \label{v:ant:eq:ii-area-enlaces}
\end{align}
Todas las coordenadas se leen módulo veintisiete. Cada colección contiene
dieciocho enlaces y ambas son disjuntas. Sus etiquetas conservan los
canales de incidencia; el cardinal de enlaces es \(18+18\).

\begin{theorem}[Capacidad mínima del corte cúbico]
\label{v:ant:thm:ii-area-corte-minimo}
En el grafo cúbico de vértices \(\{0,\ldots,N-1\}^3\), con
\(N\ge2\), enlaces entre vecinos y capacidad unitaria por enlace,
el corte mínimo entre dos caras opuestas tiene capacidad \(N^2\).
Para \(N=9\), un estado producto máximamente mezclado sobre los
trits de una capa mínima tiene información
\begin{equation}
 \operatorname{cap}_{\min}=81,\qquad s_{\rm corte}=81\log3.
 \label{v:ant:eq:ii-area-corte-informacion}
\end{equation}
\end{theorem}
\begin{proof}
Las \(N^2\) rectas paralelas a la normal de las caras forman caminos
disjuntos por enlaces. Todo corte separador intercepta cada camino, de
modo que contiene al menos \(N^2\) enlaces. Los enlaces que atraviesan
una capa transversal constituyen un corte de ese cardinal. Cada trit
uniforme aporta \(\log3\) nats; la entropía es aditiva para el estado
producto de los ochenta y un trits de la capa.
\end{proof}

La capacidad del corte cúbico y el estado de los treinta y seis enlaces
de \eqref{v:ant:eq:ii-area-enlaces} pertenecen a observables definidos sobre
cortes diferentes. La genealogía conserva ambos dominios y sus medidas.

\subsection{Transporte de los modos colectivos de incidencia}
\label{v:ant:sec:ii-area-transporte-colectivo}

Los espacios \(\mathcal F_6\) y \(\mathcal F_8\) de la
sección~\ref{v:ant:sec:barbero} proporcionan los vectores uniformes
\[
 u_{90}=\frac1{\sqrt{90}}\sum_{f\in\mathcal F_6}(\delta_f,0),
 \qquad
 u_{120}=\frac1{\sqrt{120}}\sum_{f\in\mathcal F_8}(0,\delta_f)
\]
en \(\ell^2(\mathcal F_6)\oplus\ell^2(\mathcal F_8)\).
En \(\ell^2(\partial B_{90})\oplus\ell^2(\partial B_{120})\)
se construyen
\[
 v_{90}=\frac1{\sqrt{18}}\sum_{e\in\partial B_{90}}(\delta_e,0),
 \qquad
 v_{120}=\frac1{\sqrt{18}}\sum_{e\in\partial B_{120}}(0,\delta_e).
\]
Las dos parejas son ortonormales. Denotemos por
\(\mathscr H_{\rm inc}^{\rm coll}\) y
\(\mathscr H_\partial^{\rm coll}\) sus respectivos planos generados.

\begin{theorem}[Entrelazamiento colectivo de las etiquetas sectoriales]
\label{v:ant:thm:ii-area-entrelazamiento-colectivo}
Existe una única isometría lineal sobreyectiva entre esos planos tal que
\begin{equation}
 \mathcal J_{6\to8}u_m=v_m\quad(m=90,120).
 \label{v:ant:eq:ii-area-j-colectivo}
\end{equation}
Para los operadores
\[
 \begin{aligned}
 D_{\rm inc}&=90|u_{90}\rangle\langle u_{90}|
             +120|u_{120}\rangle\langle u_{120}|,\\
 D_\partial&=90|v_{90}\rangle\langle v_{90}|
             +120|v_{120}\rangle\langle v_{120}|,
 \end{aligned}
\]
se cumple
\begin{equation}
 \mathcal J_{6\to8}D_{\rm inc}
 =D_\partial\mathcal J_{6\to8}.
 \label{v:ant:eq:ii-area-j-entrelaza}
\end{equation}
Se transportan también los proyectores sectoriales y su combinación
orientada de coeficientes \(12:-1\).
\end{theorem}
\begin{proof}
La imagen de una base fija la aplicación lineal única. Las normas y los
productos internos se conservan porque ambas bases son ortonormales.
La identidad \eqref{v:ant:eq:ii-area-j-entrelaza} vale sobre cada vector
\(u_m\), pues sus dos miembros son \(m v_m\). Los proyectores se
expresan como \((120I-D)/30\) y \((D-90I)/30\); su cálculo funcional
lineal y cualquier combinación de ambos conservan el entrelazamiento.
\end{proof}

Este transporte actúa sobre los dos modos uniformes y sus etiquetas
espectrales. Las incidencias individuales y las fluctuaciones dentro de
cada sector permanecen en sus espacios completos. El enlace conserva
exactamente el dato colectivo que utilizará la ponderación areal.

\subsection{Carácter angular y operador de área}
\label{v:ant:sec:ii-area-caracter}

Sea $N_e=\operatorname{diag}(0,1,2)$ en cada enlace y
$N_{90},N_{120}$ las sumas por sector. Fijada la escala
\[
\theta_{\rm clk}=\frac{C^*}{6}\frac{\pi}{180},\qquad
a_0=\frac{1+2\cos(10^\circ)}3\theta_{\rm clk},
\]
el factor angular es el carácter real normalizado del triplete
\begin{equation}
 \chi_{10}=\frac13\Tr\operatorname{diag}
       (1,e^{i\pi/18},e^{-i\pi/18})
 =\frac{1+2\cos(10^\circ)}3,
 \qquad a_0=\chi_{10}\theta_{\rm clk}.
 \label{v:ant:eq:ii-area-caracter-triplete}
\end{equation}
La suma de las fases conjugadas prueba la igualdad. En la rama positiva
considerada, \(\theta_{\rm clk}>0\), \(\chi_{10}>0\) y
\(\gamma_{\HMT}>0\); por tanto, \(a_0\) y la escala areal son
positivos. El triplete angular y su normalización preceden a la
distribución entrópica sobre los enlaces. Con estas cantidades,
el operador areal normalizado es
\[
\widehat{\mathcal A}/\ell_P^2
=\gamma_{\HMT}a_0(N_{90}+N_{120}).
\]
La reciprocidad de~\eqref{xii:radial:G} identifica aquí
\(\ell_P^2=L^2\) en ambas secciones de acción. Por ello la lectura
areal proporciona \(N=\widehat{\mathcal A}/(L^2\gamma_{\HMT}a_0)\);
el Hamiltoniano modular que sigue recupera la segunda coordenada.
El espectro de \(\widehat{\mathcal A}/L^2\) tiene valores $n\gamma_{\HMT}a_0$, $0\le n\le72$,
con multiplicidades dadas por $(1+z+z^2)^{36}$.
La prueba consiste en sumar los autovalores $0,1,2$ de los
36 factores tensoriales.

\begin{theorem}[Recuperación de los sectores de borde]
Sean
\[
N=N_{90}+N_{120},\qquad D=90N_{90}+120N_{120}.
\]
Entonces
\begin{equation}
N_{90}=4N-\frac D{30},\qquad
N_{120}=\frac D{30}-3N.
\label{v:ant:eq:recover}
\end{equation}
Por tanto, el total y el lector ponderado recuperan conjuntamente
las dos ocupaciones.
\end{theorem}
\begin{proof}
La matriz $\left(\begin{smallmatrix}1&1\\90&120\end{smallmatrix}\right)$
tiene determinante $30$. Su inversión da \eqref{v:ant:eq:recover}.
En el dominio de ocupaciones enteras, la imagen satisface
$D\in30\Z$ y $90N\le D\le120N$.
\end{proof}

El Hamiltoniano modular de la fuente es
$K_A=-\log\rho_A=cI+\Delta_{\rm mod}D$.
Para $\Delta_{\rm mod}\ne0$ y normalización conocida,
$D=(K_A-cI)/\Delta_{\rm mod}$ expresa el segundo observable.
Este operador describe la estructura espectral del estado reducido;
no se identifica por notación con el generador de evolución temporal.
El área conserva la ocupación total y la pareja añade la procedencia.
Por ejemplo, $(N_{90},N_{120})=(1,0)$ y $(0,1)$ comparten $N=1$,
pero tienen $D=90$ y $D=120$.

% RESULTADO_RECUPERADO desde II REV09. Véase MAPA_PROCEDENCIA.json.
% Selección de líneas y cambios mecánicos explícitos; ninguna prueba nueva.
\subsection{Normalización del estado y significado del Hamiltoniano modular}
El segundo observable se explicita mediante el estado que lo define.
Se mantiene la descomposición tensorial anterior: dieciocho qutrits
por sector, cada uno con operador número $\operatorname{diag}(0,1,2)$;
$N_m$ es la suma de esos operadores y los sectores actúan en factores
distintos. Sea
\[
 z_m=1+e^{-m\Delta_{\rm mod}}+e^{-2m\Delta_{\rm mod}},
 \qquad m\in\{90,120\}.
\]
La colección de estados de la fuente tiene la realización producto
\[
 \rho_A=Z^{-1}e^{-\Delta_{\rm mod}(90N_{90}+120N_{120})},
 \qquad Z=z_{90}^{18}z_{120}^{18}.
\]
La conmutatividad de los operadores número y la factorización tensorial
prueban la expresión de $Z$. El estado es de rango completo para
$\Delta_{\rm mod}$ real y finito, y su logaritmo espectral da
\[
 K_A=-\log\rho_A=(\log Z)I+\Delta_{\rm mod}D.
\]
Para $\Delta_{\rm mod}\ne0$, $K_A$ pesa de forma distinta las excitaciones
de los dos sectores. La lectura conjunta con $N$ recupera las dos ocupaciones
por la ecuación~\eqref{v:ant:eq:recover}.

La necesidad de la lectura conjunta se ve en dos ejemplos complementarios.
Los estados de ocupación $(1,0)$ y $(0,1)$ tienen igual total $N=1$
y diferente peso modular. A la inversa, $(4,0)$ y $(0,3)$ tienen
el mismo $D=360$ y diferentes totales. La pareja $(N,D)$ separa ambos
casos. Esta recuperación se refiere a ocupaciones sectoriales, no a
cada vector microscópico dentro de una degeneración de esas ocupaciones.

Cuando $\Delta_{\rm mod}=0$, el estado es proporcional a la identidad
y $K_A$ es escalar: esa lectura deja de distinguir los sectores.
Para $\Delta_{\rm mod}\ne0$, la normalización conocida permite recuperar
$D$ y componerlo con el área. Así queda definido el papel del Hamiltoniano
modular: caracteriza el estado reducido y su distribución de información.
Su inclusión en el artículo está motivada por esa función matemática,
mientras el Hamiltoniano de evolución pertenece al desarrollo dinámico
de la acción.


\pagebreak % REV03: keep the reduced-state explanation and equations together.
\subsection{Estado reducido, descomposición de Schmidt y entropía de entrelazamiento}
Si $\rho_A=\sum_\nu p_\nu|\nu\rangle\langle\nu|$,
la purificación
\[
|\Psi\rangle=\sum_\nu\sqrt{p_\nu}\,
|\nu\rangle_A|\nu\rangle_B
\]
satisface $\Tr_B|\Psi\rangle\langle\Psi|=\rho_A$.
Así, $-\Tr(\rho_A\log\rho_A)$ es la entropía de entrelazamiento
de esta bipartición pura concreta. El estado ampliado declara
qué sistema porta esa interpretación.

Para una involución de dos hojas con un autovector de cada signo,
\[
\rho_y=\frac{e^{-yR}}{2\cosh y},\qquad R^2=I,
\]
la diagonalización proporciona
\[
S(\rho_y)=\log(2\cosh y)-y\tanh y,\qquad
D(\rho_y\Vert\rho_{-y})=2y\tanh y.
\]
La entropía es par en $y$; la distancia compara explícitamente
las dos orientaciones. No se identifica esta entropía de dos niveles
con la capacidad $81\log3$ del corte triádico.
